% Tabla de caracterización de los dos estudios incluidos (apaisada).
% Datos extraídos del TEXTO COMPLETO (04_corpus_final/pdfs/), no del resumen.
\newgeometry{margin=2cm}
\begin{landscape}
\sloppy
\hfuzz=500pt
\emergencystretch=3em
\begin{longtblr}[
caption={\textit{Caracterización de los dos estudios incluidos, extraída del
texto completo. Las cifras corresponden a la comparación intra-estudio
reportada por los autores; no constituyen una síntesis entre estudios.}},
label={tab:studies},
]{colspec={X[1.2]X[1.7]X[3.0]X[3.0]X[3.6]X[3.3]},width=1.0\linewidth,
  hline{1,2,Z}={0.6pt},row{1}={font=\bfseries}}
Estudio & Objetivo & Estrategia de modelado y control &
Datos y variables & Resultados clave (intra-estudio) &
Validación y limitaciones \\
\cite{palaciomorales2021optimiza}
& Sintonizar un NMPC para regular el pH en la alcalinización del jugo de
  caña (\textit{setpoint} 7.34) tras identificar su modelo interno
& Modelo interno Takagi--Sugeno de entrada--salida con conjuntos difusos
  multidimensionales (consecuentes lineales); identificación y sintonía del
  NMPC (horizontes y pesos) mediante búsqueda de patrón generalizada, sin
  gradiente
& Planta \textbf{simulada} con un modelo de invariantes de reacción
  (HNO$_3$ + NaHCO$_3$) tomado de una tesis doctoral, no datos de planta.
  Manipulada: caudal de lechada de cal $F_2$ (10--70~GPM).
  Perturbaciones: caudal de jugo sulfitado $F_1$ (500--2100~GPM) y pH del
  jugo sulfitado y de la lechada. Salida: pH del jugo alcalinizado
  (5.7--8.7)
& Índice ICET(pH) $=553$ frente a $2.41$--$3.95\times10^{3}$ de otras cinco
  sintonías de NMPC publicadas (menor es mejor). Sobreimpulso
  $M_p=30.3\,\%$ frente a $29.8$--$31.1\,\%$ (sin mejora). Esfuerzo de
  control $79.2$ frente a $28.0$--$83.4$ (mayor que el de cuatro de las
  cinco). \textbf{No se compara con un PID/PI}
& Solo \textbf{simulación} (TRL~$\approx 3$--$4$, escala NASA
  NPR~7123.1B); sin planta piloto ni industrial. Los autores reportan
  rizado residual en la respuesta; el modelo de planta es un
  \textit{benchmark} genérico, no un clarificador medido \\
\cite{murugesan2020computeraided}
& Estimar los parámetros del modelo del clarificador y sintonizar fuera
  de línea un controlador PID para la neutralización del pH
& Identificación estándar (modelo de segundo orden) refinada con un
  \textit{state transition algorithm} con restricciones lagrangianas
  (LSTA); sintonía de $K_p$, $T_i$, $T_d$ minimizando el MSE; comparación
  con GA y PSO
& \textbf{Datos reales de planta} (M/S Sakthi Sugars Ltd.): 3000 muestras en
  dos turnos, 70\,\% entrenamiento y 30\,\% prueba. Entrada: caudal de
  lechada de cal. Salida: pH final del jugo
& MSE $1.2\times10^{-3}$ (LSTA) frente a $3.3\times10^{-2}$ (GA) y
  $6.6\times10^{-1}$ (PSO); tiempo de asentamiento $46.1$~s frente a $52.2$ y
  $94.8$~s; tiempo de subida $33.4$~s (LSTA) frente a $28.3$ (GA) y $37.4$~s
  (PSO), sin ventaja de LSTA en este índice
& Identificación \textbf{fuera de línea con datos de planta}; el
  controlador solo se evalúa en \textbf{simulación} (respuesta a escalón),
  sin implementación en lazo cerrado (TRL~$\approx 4$). Modelo lineal de
  segundo orden para un proceso no lineal; el controlador es PID
  sintonizado por metaheurística \\
\end{longtblr}
\nolinenumbers
\begin{Nota}[Nota]{\linewidth}{Elaboración propia a partir del texto
completo de \cite{palaciomorales2021optimiza} (DOI: 10.3390/pr9122283) y
\cite{murugesan2020computeraided} (DOI: 10.1007/s00034-019-01139-5).
Abreviaturas: NMPC $=$ \textit{nonlinear model predictive control};
ICET $=$ integral del error cuadrático ponderado por el tiempo; GA
$=$ algoritmo genético; PSO $=$ optimización por enjambre de partículas;
TRL $=$ \textit{technology readiness level}
\parencite{nasa2016trl}. Las cifras se citan tal como las reportan los
autores originales.}
\end{Nota}
\nolinenumbers
\end{landscape}
\restoregeometry
